% Part: first-order-logic
% Chapter: sequent-calculus
% Section: proving-things-quant

\documentclass[../../../include/open-logic-section]{subfiles}

\begin{document}

\olfileid{fol}{seq}{prq}

\olsection{\usetoken{P}{derivation} met kwantoren}

\begin{ex}
Geef een $\Log{LK}$-!!{derivation} van de sequent
$\lexists[x][\lnot !A(x)] \Sequent \lnot \lforall[x][!A(x)]$.

Bij kwantoren moeten we erop letten dat we de
eigenvariabelevoorwaarde niet schenden. Soms vergt dit dat we de
volgorde waarin we bepaalde afleidingsstappen uitvoeren aanpassen. In
het algemeen is het nuttig eerst regels af te handelen waarvoor de
eigenvariabelevoorwaarde geldt; in het voltooide bewijs staan die
regels dan lager. Ook is het raadzaam vooruit te kijken en te bepalen
hoe de beginsequent eruit zou kunnen zien. In ons geval moet dit iets
zijn als $!A(a) \Sequent !A(a)$. Wanneer we de kwantorregels
``omkeren'', moeten we daarom dezelfde term kiezen, die we $a$ noemen,
voor zowel de $\lforall$- als de $\lexists$-regel. Met twee
verschillende termen zouden we uitkomen op iets als $!A(a) \Sequent
!A(b)$, en die sequent is uiteraard niet afleidbaar.

Zoals gebruikelijk schrijven we eerst
\begin{prooftree}
\AxiomC{}
\UnaryInf$\lexists[x][\lnot !A(x)] \fCenter \lnot \lforall[x][!A(x)]$
\end{prooftree}
We kunnen nu de regel \LeftR{\exists} of de regel \RightR{\lnot}
toepassen. Omdat voor \LeftR{\exists} de eigenvariabelevoorwaarde
geldt, is het verstandig die regel zo vroeg mogelijk af te handelen.
Daarom passen we hem eerst toe.
\begin{prooftree}
\AxiomC{}
\UnaryInf$ \lnot !A(a) \fCenter \lnot \lforall[x][!A(x)]$
\RightLabel{\LeftR{\lexists}}
\UnaryInf$ \lexists[x][\lnot !A(x)] \fCenter \lnot \lforall[x][!A(x)]$
\end{prooftree}
Als we de regels \LeftR{\lnot} en \RightR{\lnot} achterwaarts
toepassen, verkrijgen we
\begin{prooftree}
\AxiomC{}
\UnaryInf$\lforall[x][!A(x)] \fCenter !A(a)$
\RightLabel{\LeftR{\lnot}}
\UnaryInf$\lnot !A(a), \lforall[x][!A(x)] \fCenter $
\RightLabel{\LeftR{\Exchange}}
\UnaryInf$\lforall[x][!A(x)], \lnot !A(a) \fCenter $
\RightLabel{\RightR{\lnot}}
\UnaryInf$ \lnot !A(a) \fCenter \lnot \lforall[x] !A(x)$
\RightLabel{\LeftR{\lexists}}
\UnaryInf$ \lexists[x] \lnot !A(x) \fCenter \lnot \lforall[x] !A(x)$
\end{prooftree}
Nu is onze enige mogelijkheid de regel \LeftR{\forall} toe te passen.
Voor deze regel geldt de eigenvariabelebeperking niet, zodat dit geen
probleem oplevert. We willen een beginsequent van de vorm $!A(a)
\Sequent !A(a)$ verkrijgen. Daarom kiezen we bij de toepassing van de
regel $a$ als argument van $!A$.
\begin{prooftree}
\Axiom$!A(a) \fCenter !A(a)$
\RightLabel{\LeftR{\lforall}}
\UnaryInf$\lforall[x][!A(x)] \fCenter !A(a)$
\RightLabel{\LeftR{\lnot}}
\UnaryInf$\lnot !A(a), \lforall[x][!A(x)] \fCenter $
\RightLabel{\LeftR{\Exchange}}
\UnaryInf$\lforall[x][!A(x)], \lnot !A(a) \fCenter $
\RightLabel{\RightR{\lnot}}
\UnaryInf$ \lnot !A(a) \fCenter \lnot \lforall[x][!A(x)]$
\RightLabel{\LeftR{\lexists}}
\UnaryInf$ \lexists[x][ \lnot !A(x)] \fCenter \lnot \lforall[x][!A(x)]$
\end{prooftree}
Vooral bij kwantoren is het van belang nu nogmaals te controleren dat
de eigenvariabelevoorwaarde niet is geschonden. De enige toegepaste
regel waarvoor die voorwaarde geldt, is \LeftR{\exists}. De
eigenvariabele~$a$ komt niet voor in de onderste sequent van die regel,
namelijk de eindsequent. Dit is dus een correcte !!{derivation}.
\end{ex}

\begin{prob}
Geef !!{derivation}s van de volgende sequenten:
\begin{enumerate}
\item $\Sequent (\lforall[x][!A(x)] \land \lforall[y][!B(y)]) \lif
\lforall[z][(!A(z) \land !B(z))]$.
\item $\Sequent (\lexists[x][!A(x)] \lor \lexists[y][!B(y)]) \lif
\lexists[z][(!A(z) \lor !B(z))]$.
\item $\lforall[x][(!A(x) \lif !B)] \Sequent \lexists[y][!A(y)] \lif !B$.
\item $\lforall[x][\lnot !A(x)] \Sequent \lnot\lexists[x][!A(x)]$.
\item $\Sequent \lnot\lexists[x][!A(x)] \lif \lforall[x][\lnot !A(x)]$.
\item $\Sequent \lnot\lexists[x][\lforall[y][((!A(x,y) \lif \lnot
!A(y,y)) \land (\lnot !A(y,y) \lif !A(x,y)))]]$.
\end{enumerate}
\end{prob}

\begin{prob}
Geef !!{derivation}s van de volgende sequenten:
\begin{enumerate}
\item $\Sequent \lnot\lforall[x][!A(x)] \lif \lexists[x][\lnot!A(x)]$.
\item $(\lforall[x][!A(x)] \lif !B) \Sequent \lexists[y][(!A(y) \lif !B)]$.
\item $\Sequent \lexists[x][(!A(x) \lif \lforall[y][!A(y)])]$.
\end{enumerate}
(Hiervoor is telkens de regel $\RightR{\Contraction}$ nodig.)
\end{prob}
\end{document}
